\documentclass[11pt]{amsart}
\usepackage{amsmath,amssymb,amsthm,booktabs,hyperref,url,array}
\usepackage[margin=1.05in]{geometry}
\newtheorem{theorem}{Theorem}[section]
\newtheorem{proposition}[theorem]{Proposition}
\newtheorem{lemma}[theorem]{Lemma}
\newtheorem{corollary}[theorem]{Corollary}
\newtheorem{conjecture}[theorem]{Conjecture}
\theoremstyle{definition}
\newtheorem{remark}[theorem]{Remark}
\newtheorem{definition}[theorem]{Definition}
\newcommand{\Q}{\mathbb{Q}}\newcommand{\Z}{\mathbb{Z}}\newcommand{\F}{\mathbb{F}}
\newcommand{\degphi}{\deg\varphi_E}\newcommand{\Gal}{\operatorname{Gal}}\newcommand{\ad}{\operatorname{ad}^0}
\newcommand{\Sym}{\operatorname{Sym}^2}\newcommand{\nv}{\mathrm{nv}}\newcommand{\T}{\mathbb{T}}
\title{Memorandum: the local factors behind the twist-symmetrised bound}
\author{Dharmaj Soni}
\email{dharmaj.soni@gmail.com}
\urladdr{https://orcid.org/0009-0007-6397-3351}
\date{\today}
\begin{document}
\begin{abstract}
This memorandum answers the question raised by both referee reports on the note \emph{Local component groups and the $\ell$-adic valuations of the modular degree} (since restructured and retitled \emph{Local factors of the adjoint $L$-function and the $\ell$-adic valuations of the modular degree}; the numbering below refers to the restructured version): are the twist-symmetrised local weights of its former Conjecture 1.3, now the inequality (1.3) of its Proposition 1.3, the local factors of the adjoint / congruence-ideal theory, and if not, what are the exact local factors? The answer is that the local term of a bad prime $q\ne\ell$ is the $\ell$-adic valuation of the Euler factor of the adjoint $L$-function at $s=1$ that the \emph{naive} adjoint $L$-function of Diamond, Flach and Guo omits, plus the Tamagawa exponent of the adjoint representation when the reduction is potentially multiplicative. The weights of the note were its shadow. The new rule is sharper and holds with no exception at $\ell=3,5,7,11,13$ on every optimal curve of conductor at most $500000$ and on $893$ curves outside the tables (with PARI's symmetric-square Euler factors at every bad prime, the wild primes $2$ and $3$ included). The Euler term of a prime of type $\mathrm{III}$ can be added as well except, at $\ell=3$ only, when it meets a Tamagawa term at a prime $\equiv-1\pmod3$: a configuration of a precise shape, which is exactly the one excluded by the hypotheses of Kim and Ota. Two parts of the rule are proved here: the Euler terms of the types $\mathrm{I}_0^*$ and $\mathrm{I}_n^*$, with no hypothesis beyond the irreducibility of $E[\ell]$ and $\ell\nmid 2N$ (from a formula of Diamond--Flach--Guo and an elementary twisting argument), and the Tamagawa terms for $\ell\ge5$ (from Kim--Ota). The rest is reduced to a statement about the $\Sigma$-imprimitive adjoint Selmer group that the same theory controls.
\end{abstract}
\maketitle

\section{Summary}\label{sec:summary}

\subsection{Notation}
$E/\Q$ is the $X_0(N)$-optimal curve of a non-CM isogeny class, $f=f_E$ its newform, $\ell$ an odd prime with $E[\ell]$ irreducible, $r_E$ the congruence number, so that $v_\ell(\degphi)=v_\ell(r_E)$ when $\ell^2\nmid N$ (Agashe--Ribet--Stein). $V=V_\ell E$, $\rho=\rho_{E,\ell}$, $\bar\rho$ its reduction, $A_f=\ad V$ the adjoint. We follow the conventions of Diamond, Flach and Guo \cite[\S1.7.2]{DFG04}: for a prime $p$ let $\delta_p=\dim V^{I_p}$, so $\delta_p=2$ for good reduction, $1$ for multiplicative reduction, $0$ for additive reduction; the local factor $L_p(A_f,s)$ of the adjoint $L$-function is the true one, and the \emph{naive} factor is $L_p^{\nv}(A_f,s)=L_p(A_f,s)$ if $\delta_p>0$ and $1$ if $\delta_p=0$, which is what the Euler product of $\sum a_n(f)^2n^{-s}$ produces. The \emph{exceptional} primes are those with $\delta_p=0$ and $L_p(A_f,s)\ne1$; in the normalisation of the symmetric square, $L_p(\Sym f,s)=L_p(A_f,s-1)$, and PARI's \texttt{lfunsympow(E,2)} returns these factors at every prime, including $2$ and $3$.

\begin{definition}\label{def:terms}
For a bad prime $q\ne\ell$ put
\[
e_\ell(E,q)=v_\ell\Bigl(\frac{L_q(A_f,1)^{-1}}{L_q^{\nv}(A_f,1)^{-1}}\Bigr)=v_\ell\Bigl(\frac{P_q(q^{-2})}{1-a_q(E)^2q^{-2}}\Bigr),\qquad t_\ell(E,q)=\begin{cases}v_\ell(-v_q(j_E))&\text{if }v_q(j_E)<0,\\0&\text{otherwise,}\end{cases}
\]
where $P_q(X)$ is the polynomial with $L_q(\Sym f,s)=P_q(q^{-s})^{-1}$. So $e_\ell$ is the valuation of the Euler factor that the naive $L$-function omits (it is $0$ at multiplicative primes, where the naive factor is the true one), and $t_\ell$ is the Tamagawa exponent of $\ad V$ at a potentially multiplicative prime, the $\ell$-part of the valuation of the Tate parameter, which is the index $n$ of the Kodaira symbol $\mathrm{I}_n$ or $\mathrm{I}_n^*$ at odd $q$ (at $q=2$ a symbol $\mathrm{I}_n^*$ can have potentially good reduction and then $t_\ell=0$). Both are invariant under quadratic twist.
\end{definition}

For tame primes $q\ge5$ the Euler term is read off the Kodaira type (Section~\ref{sec:euler}): with $\chi_{-3},\chi_{-4}$ the quadratic characters of $\Q(\sqrt{-3})$, $\Q(i)$, $E'$ the twist of $E$ by $q^*$, and $\alpha,\beta$ the roots of $X^2-a_q(E')X+q$ (type $\mathrm{I}_0^*$ only),
\begin{center}\footnotesize
\begin{tabular}{lll}
\toprule
type at $q\ge5$ & $L_q(A_f,s)^{-1}$ & $e_\ell(E,q)$\\
\midrule
$\mathrm{I}_n$ & $1-q^{-1-s}$ & $0$\\
$\mathrm{I}_n^*$ & $1-q^{-1-s}$ & $v_\ell(q^2-1)$\\
$\mathrm{II},\mathrm{II}^*,\mathrm{IV},\mathrm{IV}^*$ & $1-\chi_{-3}(q)q^{-s}$ & $v_\ell(q-\chi_{-3}(q))$\\
$\mathrm{III},\mathrm{III}^*$ & $1-\chi_{-4}(q)q^{-s}$ & $v_\ell(q-\chi_{-4}(q))$\\
$\mathrm{I}_0^*$ & $(1-\tfrac{\alpha}{\beta}q^{-s})(1-q^{-s})(1-\tfrac{\beta}{\alpha}q^{-s})$ & $v_\ell(q-1)+v_\ell\bigl((q+1)^2-a_q(E')^2\bigr)$\\
\bottomrule
\end{tabular}
\end{center}

\subsection{The rule}
\begin{conjecture}[Euler-factor rule]\label{conj:euler}
Let $E$ be optimal, non-CM, $\ell$ odd, $E[\ell]$ irreducible. Then
\begin{equation}\label{eq:rule}
v_\ell(\degphi)\ \ge\ \sum_{\substack{q\mid N,\ q\ne\ell\\ \text{type}\ne\mathrm{III},\mathrm{III}^*}}e_\ell(E,q)\ +\ \sum_{q\mid N}t_\ell(E,q)\ +\ [\ell=3]\,w_3(E),
\end{equation}
where $w_3(E)$ is the weight of the prime $3$ itself when $3\mid N$ and the reduction at $3$ is potentially good: $1$ for types $\mathrm{II},\mathrm{IV},\mathrm{III}^*$, $2$ for $\mathrm{II}^*,\mathrm{IV}^*$, $0$ for $\mathrm{III},\mathrm{I}_0^*$ (an empirical rule; see Section~\ref{sec:data}). The Euler term of a prime of type $\mathrm{III}$ or $\mathrm{III}^*$ can be added as well except possibly when some prime $p\equiv-1\pmod\ell$ has $t_\ell(E,p)\ge1$.
\end{conjecture}

Every term is attained as a minimum, in isolation, over tens of thousands of curves (Table~\ref{tab:iso}); \eqref{eq:rule} holds with no exception on the four ranges at $\ell=3,5,7,11,13$ (Table~\ref{tab:uniform}); with the type $\mathrm{III}$ terms included it has exceptions only at $\ell=3$ (Section~\ref{sec:anomaly}). It implies the two former conjectures of the note, its inequalities (1.2) and (1.3), because $3\mid q-\chi_{-3}(q)$ and $3\mid q^2-1$ for every $q\ge5$.

\subsection{What is proved}
Theorem~\ref{thm:twist} (twist monotonicity of the congruence number), Theorem~\ref{thm:euler} (the Euler terms of the types $\mathrm{I}_0^*$ and $\mathrm{I}_n^*$, under $\ell\nmid 2N$ and $E[\ell]$ irreducible only), Theorem~\ref{thm:C} (the Tamagawa terms, $\ell\ge5$, from Kim--Ota) and their combination, Corollary~\ref{cor:both}. For the types $\mathrm{II}$, $\mathrm{IV}$, $\mathrm{III}$ and their duals the rule is reduced in Section~\ref{sec:prop} to a statement about the $\Sigma$-imprimitive adjoint Selmer group, Proposition~\ref{prop:selmer}, which is the piece to put to an expert.

\section{The framework}\label{sec:framework}

\subsection{Diamond--Flach--Guo}
For a newform $f$ of weight $2$ and level $N$ with rational coefficients, a prime $\lambda=\ell$ not dividing $2N$ ($k!=2$), and a finite set $\Sigma$ of primes, let $N^\Sigma=N\prod_{p\in\Sigma}p^{\delta_p}$ and let $f^\Sigma=\sum_{(n,\Sigma)=1}a_nq^n$ be the $\Sigma$-depleted eigenform, of level $N^\Sigma$. The integral cohomology of $X_0(N^\Sigma)$ carries an alternating pairing; its restriction to the rank-two lattice $M_{f^\Sigma,\lambda}$ (the $f^\Sigma$-isotypic part) has a discriminant, and \cite[\S1.7.3]{DFG04} define the \emph{naive $\Sigma$-finite congruence ideal} $\eta_f^\Sigma\subset\Z_\ell$ by $\wedge^2M_{f^\Sigma,\lambda}=\eta_f^\Sigma\cdot(\text{standard})$; it is an integral ideal.

\begin{theorem}[{\cite[Proposition 1.4(c)]{DFG04}}]\label{thm:dfg14}
If $\bar\rho_f$ is irreducible then
\[
\eta^\Sigma_{f,\lambda}=\eta^{\emptyset}_{f,\lambda}\prod_{p\in\Sigma}L_p^{\nv}(A_f,1)^{-1}.
\]
\end{theorem}
Only the primes with $\delta_p\ge1$ contribute: adding an additive prime to $\Sigma$ changes neither the level nor the ideal, adding a good prime multiplies the ideal by $(1-\frac\alpha\beta q^{-1})(1-q^{-1})(1-\frac\beta\alpha q^{-1})$, Wiles's level-raising factor, and adding a multiplicative prime multiplies it by $1-q^{-2}$.

\begin{theorem}[{\cite[Theorem 0.2 $=$ 3.7]{DFG04}}]\label{thm:dfg02}
Let $\Sigma$ contain the primes dividing $N$ and let $\lambda\nmid2N\prod_{p\in\Sigma}p$ be such that $\bar\rho_f|_{\Gal(\overline\Q/\Q(\sqrt{\ell^*}))}$ is absolutely irreducible. Then the $\Sigma$-imprimitive adjoint Selmer group $H^1_\Sigma(G_\Q,A_{f,\lambda}/\mathcal{A}_{f,\lambda})$ has length $v_\lambda(\eta_f^\Sigma)$.
\end{theorem}
The same authors modify the integral structures in their Section 1.8 ``to account for the congruences corresponding to Euler factors missing from $L^{\nv}(A_f,s)$'', that is, at the exceptional primes, and prove the Bloch--Kato conjecture for $A_f$ and $A_f(1)$ (their Theorem 0.1) for $\lambda$ outside an explicit set $S_f$ (primes dividing $2N$ and primes where the irreducibility over $\Q(\sqrt{\ell^*})$ fails).

\subsection{Kim--Ota}
In the formulation of \cite[\S2]{KimOta}, which follows \cite{DFG04} and Dimitrov: under Assumption 2.1 ((FL): $k\le\ell-1$ and $\ell\nmid N$; (TW): $\bar\rho|_{\Gal(\overline\Q/\Q(\sqrt{\ell^*}))}$ absolutely irreducible) and for $\Sigma$ containing their exceptional set $P_{\bar\rho}$, the $\Sigma$-ramified deformation ring is the Hecke algebra $T_\Sigma$, a complete intersection, the Hecke module is free of rank two over it, and for every newform $f$ in $T_\Sigma$, $\#\mathcal{O}/\eta_{f,\Sigma}=\#\mathrm{Sel}_\Sigma(\Q,\ad f\otimes E/\mathcal{O})$ (their Theorem 2.9). Their Lemma 2.7 identifies the congruence ideal of the depleted form $f^\Sigma$ at level $N_f^\Sigma$ in the full Hecke algebra with $\eta_{f,\Sigma}$. Their main theorem (1.1) is the quantitative level lowering recalled in the note: for $N=N^+N^-$ with $N^-$ square-free, $\ell\ge5$, $\ell\nmid N$, (TW), and \emph{(6)}: $\bar\rho$ ramified at every $q\mid N^-$ with $q\equiv\pm1\pmod\ell$,
\begin{equation}\label{eq:ko}
\operatorname{ord}_\lambda\eta_f(N)=\operatorname{ord}_\lambda\eta_f(N^+,N^-)+\sum_{q\mid N^-}t_f(q),
\end{equation}
with $\eta_f(N)$ the congruence ideal in the full space and $t_f(q)$ the Tamagawa exponent (Definition 2.12 there), which is $t_\ell(E,q)$ for an elliptic curve. Remark 2.2 there, after \cite[\S1.7.1]{DFG04}, says that one may assume the newforms to have minimal conductor among their twists when working with $\ad f$: this is the structural reason for the twist-symmetrised weights of the note.

\subsection{Darmon--Diamond--Taylor}
Three statements from \cite{DDT} are used. (i) \emph{The algebra of Lemma 4.17}: let $\T$ be a Hecke algebra acting on a lattice $L$ free of finite rank over $\mathcal{O}$ with a perfect pairing for which $\T$ is self-adjoint, $\pi\colon\T\to\mathcal{O}$ an eigencharacter, $\wp=\ker\pi$, $I=\operatorname{Ann}_\T(\wp)$, $\eta=\pi(I)$ the congruence ideal, and suppose $L[\wp]$ is free of rank $d$ with Gram determinant $\Delta$. Then $\eta^d\subseteq\Delta\mathcal{O}$, with equality if $L$ is free over $\T$. The proof is three lines: $L/(L[\wp]\oplus L[I])$ is killed by $\eta$, generated by $d$ elements, and of order $\#\mathcal{O}/\Delta$ by the duality $L/L[I]\simeq\operatorname{Hom}(L[\wp],\mathcal{O})$. (ii) \emph{The Greenberg--Wiles formula}, Theorem 2.18: for local conditions $\mathcal{L}$ on a finite module $M$, $\#H^1_{\mathcal{L}}(\Q,M)/\#H^1_{\mathcal{L}^*}(\Q,M^*)=\#H^0(\Q,M)/\#H^0(\Q,M^*)\cdot\prod_v\#\mathcal{L}_v/\#H^0(G_v,M)$. (iii) Corollary 2.43: the tangent space of $R_\Sigma$ has dimension $\dim H^1_\Sigma(\Q,\ad\bar\rho(1))+d_\ell+\sum_{p\in\Sigma-\{\ell\}}\dim H^0(\Q_p,\ad\bar\rho(1))$, ``by the local Euler characteristic formula''. The terms $H^0(\Q_p,\ad\bar\rho(1))$ are the residual versions of the Euler terms of Definition~\ref{def:terms}: for the $\ell$-adic module, $\#H^0(\Q_p,\ad\rho(1)\otimes K/\mathcal{O})=\#\mathcal{O}/\det(1-\mathrm{Frob}_p\,|\,(\ad V)^{I_p}(1))$ when finite, and that determinant is the inverse Euler factor of $A_f$ at $s=1$. The Chapter 4 of \cite{DDT} in which Theorem 4.20 (Hida's formula) is proved assumes $\#\bar\rho(I_p)\mid\ell$ at every $p\ne\ell$ (their Section 3.3, hypothesis (e)); that excludes most of our additive types, which is why the proofs below use \cite{DFG04} instead.

\section{The local terms}\label{sec:euler}

\subsection{Tame primes}
Let $q\ge5$, $q\ne\ell$. The inertia group acts on $V$ through a finite cyclic quotient for potentially good reduction, and through a quadratic character times a unipotent for potentially multiplicative reduction. The invariants $(\ad V)^{I_q}$ and the Frobenius eigenvalue $\lambda_q$ on them, in the arithmetic normalisation in which a good prime has eigenvalues $\alpha/\beta,1,\beta/\alpha$ on $\ad V$, give $L_q(A_f,s)^{-1}=\prod(1-\lambda q^{-s})$; the values of the table in Section~\ref{sec:summary} are obtained as follows, and were confirmed against PARI's \texttt{lfunsympow} on every curve of the tables.

\emph{Multiplicative, $\mathrm{I}_n$.} The local component of $f$ is an unramified twist of the Steinberg representation; $L_q(A_f,s)=(1-q^{-1-s})^{-1}$, which is also the naive factor since $a_q^2=1$. So $e_\ell=0$. \emph{$\mathrm{I}_n^*$.} $V=V'\otimes\chi_q$ with $V'$ Steinberg at $q$, $\ad V=\ad V'$, $a_q=0$: $e_\ell=v_\ell(1-q^{-2})=v_\ell(q^2-1)$. \emph{$\mathrm{II},\mathrm{II}^*,\mathrm{IV},\mathrm{IV}^*$.} The image of $I_q$ in $\operatorname{Aut}V$ is cyclic of order $3$ or $6$, generated by $\sigma$ or $-\sigma$ with $\sigma$ of order $3$, eigenvalues $\zeta_3^{\pm1}$. On $\ad V$ the sign disappears and the invariants form the line spanned by $\sqrt{-3}\in\Z_\ell[\zeta_3]\simeq\operatorname{End}_{I_q}(V)$; Frobenius centralises $\sigma$ if $q\equiv1\pmod3$ and inverts it if $q\equiv2$, acting on the line by $\chi_{-3}(q)$. So $e_\ell=v_\ell(1-\chi_{-3}(q)q^{-1})=v_\ell(q-\chi_{-3}(q))$. \emph{$\mathrm{III},\mathrm{III}^*$.} Image of order $4$, eigenvalues $\pm i$; the invariant line of $\ad V$ is the diagonal element, Frobenius acts by $\chi_{-4}(q)$; $e_\ell=v_\ell(q-\chi_{-4}(q))$. \emph{$\mathrm{I}_0^*$.} $V=V'\otimes\chi_q$ with $V'$ unramified, Frobenius eigenvalues $\alpha,\beta$; $\ad V$ unramified with eigenvalues $\alpha/\beta,1,\beta/\alpha$, so $L_q(A_f,1)^{-1}=(1-\frac{\alpha}{\beta q})(1-\frac1q)(1-\frac{\beta}{\alpha q})=\frac{q-1}{q}\cdot\frac{(q+1)^2-a_q(E')^2}{q^2}$.

\subsection{Wild primes and the prime $\ell$}
At $q=2,3$ the Kodaira type does not determine $L_q(A_f,s)$, but PARI does, and Definition~\ref{def:terms} applies verbatim; for instance a type $\mathrm{IV}$ at $2$ has $P_2(X)=1+2X$ and $e_3=v_3(1+\tfrac12)=1$, a type $\mathrm{II}$ at $2$ has $P_2=1$ and $e_3=0$, and a type $\mathrm{I}_n^*$ at $2$ has potentially good reduction in about half of the cases, with a cubic $P_2$ and $t=0$ (this resolves the irregular behaviour of the prime $2$ reported in the note). At $q=\ell$ the Euler factor is not the relevant quantity (the term at $\ell$ in Hida's formula is different), and the weights $w_3$ of Conjecture~\ref{conj:euler} are empirical.

\section{Twist monotonicity}\label{sec:twist}

\begin{theorem}\label{thm:twist}
Let $E/\Q$ have conductor $N$, $S$ a set of odd primes, $D=\prod_{q\in S}q$, $\chi_S$ the product of the quadratic characters of conductor $q\in S$, $E_S=E\otimes\chi_S$ of conductor $N_S$. If $q^2\mid N$ for every $q\in S$ and $\operatorname{lcm}(N_S,D^2)\mid N$, then $r_{E_S}\mid r_E$; if moreover $\operatorname{lcm}(N,D^2)\mid N_S$, then $r_{E_S}=r_E$.
\end{theorem}
\begin{proof}
Let $f=f_E$, $g=f_{E_S}$ and $\iota\colon h\mapsto h\otimes\chi_S=\sum\chi_S(n)a_n(h)q^n$, which maps $S_2(\Gamma_0(N_S))$ into $S_2(\Gamma_0(\operatorname{lcm}(N_S,D^2)))\subset S_2(\Gamma_0(N))$, preserves integrality, and sends eigenforms to eigenforms with eigenvalues twisted by $\chi_S(p)$ at $p\nmid D$ and coefficients at multiples of $D$ killed. Then $\iota(g)=f$: both are multiplicative, $a_p(f)=\chi_S(p)a_p(g)$ for $p\nmid D$ since the twist is unramified outside $S$, and $a_{q^k}(f)=0$ for $q\in S$ since $q^2\mid N$. An eigenform $h\ne g$ of level dividing $N_S$ has an eigensystem different from that of $g$ at infinitely many primes, so $\iota(h)\in f^\perp$. Hence a congruence $g\equiv h\pmod r$ with $h\in g^\perp$ integral gives $f\equiv\iota(h)\pmod r$ with $\iota(h)\in f^\perp$ integral, and $r\mid r_E$ by the cyclicity of the congruence module (Lemma 3.1 of the note). With $r=r_{E_S}$ this is the first claim; the second hypothesis makes the situation symmetric.
\end{proof}
The hypotheses hold when (a) $E$ has type $\mathrm{I}_n^*$ at each $q\in S$ ($N_S=N/D$), (b) type $\mathrm{I}_0^*$ at each $q\in S$ ($N_S=N/D^2$, $E_S$ good at $S$), (c) type $\mathrm{II}$, $\mathrm{III}$, $\mathrm{IV}$ or a dual at each $q\in S$, $q\ge5$ ($N_S=N$, equality). On the working set, for the $157{,}938$ pairs (curve, prime of type $\mathrm{I}_n^*$) and the $292{,}580$ pairs (curve, prime of potentially good additive type), $v_\ell(\degphi)\ge v_\ell(\deg\varphi_{E_q})$ holds for $\ell\in\{3,5,7,11,13\}$, $\ell^2\nmid N$, with no exception, with equality in every pair of case (c) and strict inequality only in case (b), where the difference is the Euler term of Theorem~\ref{thm:euler} (script \texttt{step9\_twist\_degree.py}).

\section{The Euler terms of the types \texorpdfstring{$\mathrm{I}_0^*$ and $\mathrm{I}_n^*$}{I0* and In*}}\label{sec:proved}

\begin{theorem}\label{thm:euler}
Let $\ell$ be an odd prime with $\ell\nmid 2N$ and $E[\ell]$ irreducible. Let $S$ be a set of primes $q\ge5$ at which $E$ has type $\mathrm{I}_0^*$ or $\mathrm{I}_n^*$, and $E_S=E\otimes\chi_S$. Then
\[
v_\ell(\degphi)\ \ge\ v_\ell(r_{E_S})\ +\ \sum_{q\in S}e_\ell(E,q),
\]
with $e_\ell(E,q)=v_\ell(q-1)+v_\ell((q+1)^2-a_q(E_S)^2)$ for type $\mathrm{I}_0^*$ and $v_\ell(q^2-1)$ for type $\mathrm{I}_n^*$.
\end{theorem}

\begin{proof}
Write $f=f_E$, $g=f_{E_S}$ of level $N_S$, $D=\prod_{q\in S}q$, and $\Sigma=S$. By Lemma 2.4 of the note, $E_S$ has good reduction at the $\mathrm{I}_0^*$ primes of $S$ and multiplicative reduction at the $\mathrm{I}_n^*$ primes, so $\delta_q=2$, resp.\ $1$, for $g$ at $q\in S$, and $N_S^\Sigma=N_S\prod_{q\in S}q^{\delta_q}=N$: the $S$-depleted form $g^S$ lives at level $N$, and $\iota(g^S)=f$ with $\iota$ the twist map of Theorem~\ref{thm:twist} (its proof only used $a_{q^k}(f)=0$ for $q\in S$ and $\iota$ on the prime-to-$D$ coefficients, so it applies to $g^S$ as well as to $g$).

\emph{Step 1: the twist map and the full Hecke algebra.} Let $\T$ be the full Hecke algebra of level $N$ and $\mathfrak{m}$ the maximal ideal of $g^S$, which contains $U_q$ for $q\in S$. A normalised eigenform for $\T$ whose eigencharacter is congruent to that of $g^S$ modulo $\mathfrak m$ has $U_q\equiv0$ for $q\in S$. The $q$-stabilisations of $g$ other than $g^S$ have $U_q$-eigenvalue $a_q(g)=\pm1$ (multiplicative case) or a root $\alpha,\beta$ of $X^2-a_q(g)X+q$ (good case), which are $\ell$-adic units since $\ell\ne q$; so they are not congruent to $g^S$ modulo $\mathfrak m$, and the forms congruent to $g^S$ are either new at every $q\in S$ with $q^2$ in their level, or old from newforms $h\neq g$ at primes outside $S$. For both kinds $\iota(h)$ is an eigenform of level dividing $N$ with eigensystem different from that of $f$ (at $p\nmid DN$ the eigensystems of $h$ and $g$ differ at infinitely many primes; at $q\in S$ both have $U_q=0$), hence $\iota(h)\in f^\perp$. Therefore the argument of Theorem~\ref{thm:twist} gives: if $g^S\equiv h\pmod{\lambda^t}$ coefficientwise for an integral $h$ in the full-Hecke-algebra orthogonal complement of $g^S$ supported on the eigenforms above, then $f\equiv\iota(h)\pmod{\lambda^t}$ with $\iota(h)\in f^\perp$ integral. By the duality $S_2(\Gamma_0(N);\Z)\simeq\operatorname{Hom}(\T,\Z)$ and the cyclicity argument of Lemma 3.1 of the note (applied to the eigencharacter $\pi\colon\T\to\Z$ of $g^S$, which has $\T/\ker\pi=\Z$), the $\ell$-part of the congruence number of $g^S$ in $\T$, namely $v_\ell(\pi(\operatorname{Ann}_\T\ker\pi))=:v_\ell(\eta^{\mathrm{Hecke}}(g^S))$, is attained by such $h$, and $v_\ell(r_E)\ge v_\ell(\eta^{\mathrm{Hecke}}(g^S))$.

\emph{Step 2: Hecke versus cohomological congruence ideal.} Apply the algebra of \cite[Lemma 4.17]{DDT} to $L=H_1(X_0(N),\Z)\otimes\mathcal{O}$ localised at $\mathfrak m$, with its intersection pairing made $\T$-self-adjoint as in \cite[\S4.4]{DDT}, and to $\pi$: $L[\wp]$ is the rank-two lattice $M_{g^S,\lambda}$ of \cite[\S1.7.3]{DFG04}, whose Gram determinant is $(\eta^S_{g,\lambda})^2$ up to a unit, the pairing being alternating. Hence $\eta^{\mathrm{Hecke}}(g^S)^2\subseteq(\eta^S_{g,\lambda})^2$, that is, $v_\ell(\eta^{\mathrm{Hecke}}(g^S))\ge v_\lambda(\eta^S_{g,\lambda})$. (Freeness is not needed for this inequality.)

\emph{Step 3: the formula of Diamond, Flach and Guo.} By Theorem~\ref{thm:dfg14}, applied to $g$ ($\bar\rho_g=\bar\rho\otimes\chi_S$ is irreducible, $\ell\nmid 2N_S$),
\[
v_\lambda(\eta^S_{g,\lambda})=v_\lambda(\eta^\emptyset_{g,\lambda})+\sum_{q\in S}v_\ell\bigl(L_q^{\nv}(A_g,1)^{-1}\bigr)=v_\lambda(\eta^\emptyset_{g,\lambda})+\sum_{q\in S}e_\ell(E,q),
\]
the last equality because $A_g=A_f$ and, at $q\in S$, $g$ has $\delta_q\ge1$ so that its naive factor is the true one, which is the factor that $f$ (with $a_q(f)=0$) omits.

\emph{Step 4: the level $N_S$.} At level $N_S$ the form $g$ is new, $\ell\nmid N_S$ and $\bar\rho_g$ is irreducible with Serre weight $2\ne\ell$, so multiplicity one holds for the maximal ideal of $g$ (\cite[Theorem 3.5]{RibetStein} and the references in its proof: the only possible exception there has Serre weight $\ell$, not $2$), hence the localised Hecke algebra is Gorenstein and $H_1(X_0(N_S),\mathcal{O})_{\mathfrak m}$ is free of rank two over it; the equality case of Lemma 4.17 then gives $v_\lambda(\eta^\emptyset_{g,\lambda})=v_\ell(r_{E_S})$. Putting the steps together and using $v_\ell(\degphi)=v_\ell(r_E)$ (as $\ell\nmid N$) proves the theorem. Without Step 4 one still has $v_\ell(\degphi)\ge\sum_{q\in S}e_\ell(E,q)$, since $\eta^\emptyset_{g,\lambda}$ is an integral ideal.
\end{proof}

\begin{theorem}\label{thm:C}
Let $\ell\ge5$, $\ell\nmid N$, with $\bar\rho_{E,\ell}$ absolutely irreducible on $\Gal(\overline\Q/\Q(\sqrt{\ell^*}))$. Let $M$ be the set of primes $q$ at which $E$ has type $\mathrm{I}_n$, or $\mathrm{I}_n^*$ with $q$ odd, with $\ell\mid n$, and $M'\subset M$ the subset of primes $q\not\equiv\pm1\pmod\ell$. Then $v_\ell(r_{E_S})\ge\sum_{q\in M'}t_\ell(E,q)$ for $S$ the set of $\mathrm{I}_n^*$ primes in $M'$, and hence $v_\ell(\degphi)\ge\sum_{q\in M'}t_\ell(E,q)$.
\end{theorem}
\begin{proof}
$E_S$ is multiplicative at every prime of $M'$ with the same index, has conductor prime to $\ell$, and $\bar\rho_{E_S}=\bar\rho\otimes\chi_S$ satisfies (TW). Apply \eqref{eq:ko} to $f_{E_S}$ with $N^-=\prod_{q\in M'}q$: square-free, prime to $N^+$, and (6) holds vacuously. Then use Theorem~\ref{thm:twist}(a).
\end{proof}

\begin{corollary}\label{cor:both}
Under the hypotheses of Theorem~\ref{thm:C}, let $S$ be the set of primes $q\ge5$ of type $\mathrm{I}_0^*$ or $\mathrm{I}_n^*$. Then
\[
v_\ell(\degphi)\ \ge\ \sum_{q\in S}e_\ell(E,q)+\sum_{q\in M'}t_\ell(E,q):
\]
the Euler term and the Tamagawa term of a prime of type $\mathrm{I}_n^*$ add. For $\ell=3$ and any set $S$ of primes $q\ge5$ of type $\mathrm{I}_n^*$ or $\mathrm{I}_0^*$, Theorem~\ref{thm:euler} alone gives $v_3(\degphi)\ge\sum_{q\in S}e_3(E,q)$, so a single prime of type $\mathrm{I}_n^*$, with $3\mid n$ or not, forces $3\mid\degphi$, and a prime of type $\mathrm{I}_0^*$ with $q\equiv1\pmod3$ does too.
\end{corollary}
\begin{proof}
Theorem~\ref{thm:euler} for $S$ gives $v_\ell(\degphi)\ge v_\ell(r_{E_S})+\sum_S e_\ell$; the twist $E_S$ is multiplicative at $S$ and Theorem~\ref{thm:C} applied to $E_S$ (whose $\mathrm{I}_n$ primes include $M'$ with the same indices) gives $v_\ell(r_{E_S})\ge\sum_{M'}t_\ell$, where for the $\mathrm{I}_n^*$ primes of $M'$ no further twist is needed.
\end{proof}

\begin{remark}
(1) Theorem~\ref{thm:euler} is the quantitative form of the twist-and-lower argument of the note, read in the other direction: the level is raised, from $N_S$ to $N_S\prod q^{\delta_q}=N$, and the price is the level-raising Euler factor. (2) The only inputs are the formula of \cite{DFG04} for the naive congruence ideal, which needs $\bar\rho$ irreducible and $\ell\nmid 2N$, the algebraic inequality of \cite[Lemma 4.17]{DDT}, and multiplicity one at the lower level. In particular $\ell=3$ is allowed, and no Taylor--Wiles hypothesis enters. (3) Step 1 is where the argument would fail for the potentially good types $\mathrm{II}$, $\mathrm{IV}$, $\mathrm{III}$: there the twist does not lower the level, $f$ is its own depletion, and Theorem~\ref{thm:dfg14} gives no factor; the Euler factor sits inside the naive $L$-value, which is where Section~\ref{sec:prop} picks it up.
\end{remark}

\section{The remaining types: a Selmer-group proposition}\label{sec:prop}

Let $\Sigma$ be the set of primes dividing $N$ and $M=\ad V\otimes K/\mathcal{O}$. Under the hypotheses of Theorem~\ref{thm:dfg02}, $v_\ell(\eta_f^\Sigma)=\operatorname{length}H^1_\Sigma(\Q,M)$, and by Theorem~\ref{thm:dfg14}, $v_\ell(\eta_f^\Sigma)=v_\ell(r_E)+\sum_{p\,\|\,N}v_\ell(p^2-1)$ (multiplicity one at level $N$ identifying $\eta^\emptyset_f$ with $r_E$). So the Euler-factor rule, for all types at once, is the statement
\begin{equation}\label{eq:selmer}
\operatorname{length}H^1_\Sigma(\Q,M)\ \ge\ \sum_{p\in\Sigma,\ \delta_p=0}e_\ell(E,p)\ +\ \sum_{p\in\Sigma,\ \delta_p=1}\bigl(t_\ell(E,p)+v_\ell(p^2-1)\bigr).
\end{equation}
Compare $H^1_\Sigma$ with the Bloch--Kato Selmer group $H^1_f(\Q,M)$ through the Greenberg--Wiles formula applied to the finite layers: $H^1_\Sigma\supseteq H^1_f$, the local index at $p\in\Sigma$ between the unrestricted condition and the Bloch--Kato (unramified) condition is, by local Tate duality and the local Euler characteristic formula, $\#H^0(\Q_p,M^*)=\#H^0(\Q_p,\ad\rho(1)\otimes K/\mathcal{O})$ when the latter is finite, which is the $\ell$-part of $L_p(A_f,1)^{-1}$, i.e.\ $\ell^{e_\ell(E,p)}$ at an exceptional prime and $\ell^{v_\ell(p^2-1)+\dots}$ at a Steinberg prime; and the global correction is the ratio of the dual Selmer groups $\#H^1_{f^*}(\Q,M^*)/\#H^1_{\Sigma^*}(\Q,M^*)\ge1$, with $M^*=\ad\rho(1)\otimes K/\mathcal{O}$ the module whose Bloch--Kato conjecture \cite{DFG04} also prove. Hence:

\begin{proposition}[the statement to prove]\label{prop:selmer}
With $\ell\nmid2N$ and (TW): the Euler-factor rule for the types $\mathrm{II}$, $\mathrm{IV}$, $\mathrm{III}$ and their duals is equivalent to
\begin{multline*}
\operatorname{length}H^1_f(\Q,\ad\rho\otimes K/\mathcal{O})\ +\ \sum_{p\in\Sigma}\operatorname{length}H^0(\Q_p,\ad\rho(1)\otimes K/\mathcal{O})\\
-\ \operatorname{length}\frac{H^1_{f^*}(\Q,\ad\rho(1)\otimes K/\mathcal{O})}{H^1_{\Sigma^*}(\Q,\ad\rho(1)\otimes K/\mathcal{O})}\ \ge\ \text{right-hand side of \eqref{eq:selmer}},
\end{multline*}
and the data of Section~\ref{sec:data} say that the Bloch--Kato Selmer group supplies the Tamagawa exponents and the dual correction vanishes, except in the configuration of Section~\ref{sec:anomaly}, where it equals one.
\end{proposition}

This is the precise form of the referees' question: the twist-symmetrised weights are explained by the twist invariance of $\ad\rho$, the exact local factors are the $H^0(\Q_p,\ad\rho(1)\otimes K/\mathcal{O})$ of Corollary 2.43 of \cite{DDT} (whose residual version is the Taylor--Wiles count), and the one thing not in the literature as far as we can see is the vanishing of the dual correction, whose failure the type $\mathrm{III}$ anomaly exhibits. The refined integral structures of \cite[\S1.8]{DFG04}, built precisely to account for the missing Euler factors at exceptional primes, are the natural place to settle it.

\section{The data}\label{sec:data}

All computations are on the working set ($N\le300000$), hold-out ($300000<N\le400000$), extension ($400000<N\le500000$) and the $893$ out-of-table curves of the note, with $E$ optimal, non-CM, $E[\ell]$ irreducible; scripts \texttt{step9\_*.py} (Kodaira-type formulas, $q\ge5$) and \texttt{step10\_*.py} (PARI's Euler factors at every bad prime) in the repository, outputs in \texttt{results/}.

\subsection{Each term in isolation}
For curves with exactly one additive prime $q\ge5$, $q\ne\ell$, and no multiplicative prime with $\ell\mid n$, Table~\ref{tab:iso} gives the minimum of $v_\ell(\degphi)$ by type and by $e_\ell(E,q)$. In every cell, for $\ell=3,5,7$, the minimum equals $e_\ell$ and the proportion of curves with $v_\ell(\degphi)\ge e_\ell$ is $1$. For $\mathrm{I}_n^*$ the excess over $e_\ell$ has minimum exactly $v_\ell(n)$ in every cell. Multiplicative primes show no Euler term: for curves with a single multiplicative prime $q$ with $\ell\mid n$ and no additive prime, $\min(v_\ell(\degphi)-v_\ell(n))=0$ for every value of $v_\ell(q^2-1)$ up to $6$.

\begin{table}[ht]\centering\footnotesize
\begin{tabular}{l rr rr rr rr}
\toprule
 & \multicolumn{2}{c}{$e=1$} & \multicolumn{2}{c}{$e=2$} & \multicolumn{2}{c}{$e=3$} & \multicolumn{2}{c}{$e=4$}\\
type ($\ell=3$) & min & $n$ & min & $n$ & min & $n$ & min & $n$\\
\midrule
$\mathrm{II}$ & 1 & 8,327 & 2 & 821 & 3 & 59 & 4 & 2\\
$\mathrm{II}^*$ & 1 & 4,184 & 2 & 296 & 3 & 12 & 4 & 1\\
$\mathrm{IV}$ & 1 & 4,290 & 2 & 296 & 3 & 12 & 4 & 1\\
$\mathrm{IV}^*$ & 1 & 8,221 & 2 & 821 & 3 & 59 & 4 & 2\\
$\mathrm{III}$ & 1 & 1,271 & 2 & 76 & 3 & 9 & -- & --\\
$\mathrm{III}^*$ & 1 & 1,271 & 2 & 76 & 3 & 9 & -- & --\\
$\mathrm{I}_0^*$ & 1 & 1,370 & 2 & 4,432 & 3 & 2,737 & 4 & 247\\
$\mathrm{I}_n^*$ & 1 & 30,497 & 2 & 2,589 & 3 & 145 & 4 & 13\\
\bottomrule
\end{tabular}
\caption{Working set, $\ell=3$: minimum of $v_3(\degphi)$ (number of curves) over curves whose only additive prime is a single $q\ge5$ of the given type with $e_3(E,q)=e$ and no multiplicative $3$-contribution. The cells with $e=0$ ($\mathrm{III}$: 5,534 curves, $\mathrm{I}_0^*$: 5,394) have minimum $0$; for $\mathrm{I}_0^*$ the term runs up to $e=6$ (minimum $6$). The same holds at $\ell=5$ and $7$ (\texttt{results/step9\_euler\_l5.txt}, \texttt{l7}).}\label{tab:iso}
\end{table}

\subsection{The rule on whole ranges}
Table~\ref{tab:uniform} checks \eqref{eq:rule} with PARI's Euler factors at every bad prime $q\ne\ell$ and the terms $t_\ell$, $w_3$ as defined, on the four ranges: no violation at any $\ell$; and it reports what happens when the type $\mathrm{III}$ terms are added (exceptions at $\ell=3$ only) and when the weights at $q=\ell$ are dropped (none needed at $\ell\ge5$). The equality rate at $\ell=3$ is $57\%$ on the working set, against $39\%$ for the twist rule of the note and $48\%$ for the version of \eqref{eq:rule} restricted to $q\ge5$.

\begin{table}[ht]\centering\footnotesize
\begin{tabular}{llrrrrrr}
\toprule
set & $\ell$ & curves & \eqref{eq:rule} $>0$ & viol. & equality & viol.\ $+\mathrm{III}$ & no $q=\ell$: $>0$ / viol. / equality\\
\midrule
working set & 3 & 1,233,729 & 861,074 & 0 & 56.8\% & 364 & 800,833 / 0 / 52.1\% \\
 & 5 & 1,303,224 & 330,477 & 0 & 75.5\% & 0 & 330,477 / 0 / 75.5\% \\
 & 7 & 1,307,753 & 173,488 & 0 & 84.1\% & 0 & 173,488 / 0 / 84.1\% \\
 & 11 & 1,309,103 & 59,495 & 0 & 91.5\% & 0 & 59,495 / 0 / 91.5\% \\
 & 13 & 1,308,945 & 35,635 & 0 & 92.4\% & 0 & 35,635 / 0 / 92.4\% \\
\midrule
hold-out & 3 & 409,577 & 295,973 & 0 & 53.9\% & 68 & 277,265 / 0 / 49.3\% \\
 & 5 & 428,251 & 114,860 & 0 & 74.0\% & 0 & 114,860 / 0 / 74.0\% \\
 & 7 & 429,280 & 60,580 & 0 & 83.2\% & 0 & 60,580 / 0 / 83.2\% \\
 & 11 & 429,502 & 21,097 & 0 & 90.6\% & 0 & 21,097 / 0 / 90.6\% \\
 & 13 & 429,498 & 12,681 & 0 & 92.3\% & 0 & 12,681 / 0 / 92.3\% \\
\midrule
extension & 3 & 348,227 & 255,507 & 0 & 53.2\% & 43 & 238,359 / 0 / 48.1\% \\
 & 5 & 358,244 & 99,401 & 0 & 73.2\% & 0 & 99,401 / 0 / 73.2\% \\
 & 7 & 358,717 & 52,766 & 0 & 82.7\% & 0 & 52,766 / 0 / 82.7\% \\
 & 11 & 358,830 & 19,136 & 0 & 90.6\% & 0 & 19,136 / 0 / 90.6\% \\
 & 13 & 358,816 & 11,484 & 0 & 91.9\% & 0 & 11,484 / 0 / 91.9\% \\
\midrule
out of table & 3 & 893 & 827 & 0 & 51.6\% & 0 & 804 / 0 / 42.3\% \\
 & 5 & 893 & 111 & 0 & 82.0\% & 0 & 111 / 0 / 82.0\% \\
 & 7 & 893 & 52 & 0 & 90.4\% & 0 & 52 / 0 / 90.4\% \\
 & 11 & 893 & 19 & 0 & 94.7\% & 0 & 19 / 0 / 94.7\% \\
 & 13 & 893 & 12 & 0 & 100.0\% & 0 & 12 / 0 / 100.0\% \\
\bottomrule
\end{tabular}

\caption{The Euler-factor rule with PARI's factors at all primes: curves with positive bound, violations and equality rate for \eqref{eq:rule}; violations when the $\mathrm{III}$ terms are added at every prime (``$+\mathrm{III}$''); violations and equality rate without the weights at $q=\ell$ (``no $q=\ell$'').}\label{tab:uniform}
\end{table}

\subsection{The type \texorpdfstring{$\mathrm{III}$}{III} anomaly}\label{sec:anomaly}
With the $\mathrm{III}$ terms included, all exceptions are at $\ell=3$ and have one shape: a prime $q$ of type $\mathrm{III}$ or $\mathrm{III}^*$ whose Euler factor is nontrivial with Frobenius acting trivially on the invariant line ($q\equiv1\pmod{12}$ for $q\ge5$; at $q=2$ the factor $1+2X$), together with a prime $p\equiv-1\pmod3$ with $t_3(E,p)\ge1$ (a multiplicative prime, or one of type $\mathrm{I}_n^*$ with $3\mid n$), and $v_3(\degphi)$ falls short by exactly one. On the clean configurations of the working set (one such $\mathrm{III}$ prime, one such $p$, nothing else; 1,066 curves) the failure rate is $2\%$, it is $0$ for $q\equiv11\pmod{12}$ (646 curves), and it rises with $v_3(q-1)$ and with $v_3(p+1)$ (to $12\%$ for $p\equiv8\pmod9$); it does not depend on whether the multiplicative reduction at $p$ is split. For every other type the Euler term adds to the Tamagawa terms with no exception, in particular when all the Tamagawa primes are $\equiv-1\pmod3$ (14,000 curves). Type $\mathrm{III}$ is the only type whose residual inertia image mod $3$ is semisimple and non-scalar; in the language of Proposition~\ref{prop:selmer}, the exceptions are the cases where the dual Selmer correction is $1$.

\subsection{The prime \texorpdfstring{$\ell$}{l} and the wild primes}
For $\ell=5,7$ no weight is needed at $q=\ell$ or at $q=2,3$ beyond Definition~\ref{def:terms}: the minima over curves with a single wild additive prime are $0$ except $t_\ell$ for potentially multiplicative reduction at $3$. For $\ell=3$ and $q=3$ (potentially good) the minima are $1$ for $\mathrm{II},\mathrm{IV},\mathrm{III}^*$, $2$ for $\mathrm{II}^*,\mathrm{IV}^*$ and $0$ for $\mathrm{III},\mathrm{I}_0^*$, which are the weights $w_3$; at $q=2$ Definition~\ref{def:terms} with PARI's factor reproduces every minimum of the note (type $\mathrm{IV}$ at $2$: $e_3=1$; type $\mathrm{II}$ at $2$: $0$; type $\mathrm{I}_n^*$ at $2$: $t_3=v_3(-v_2(j))$ when potentially multiplicative, else $0$, which explains why the index $n$ of the symbol was not the right invariant).

\section{Comparison table}\label{sec:table}
Table~\ref{tab:cmp} sets the weights of the note against the exact local terms and records what is proved for each type.
\begin{table}[htbp]\centering\footnotesize
\begin{tabular}{p{2.1cm}p{1.6cm}p{2.5cm}p{3.1cm}p{3.8cm}}
\toprule
type at $q\ge5$ & weight in the note & local representation & exact local term & status\\
\midrule
$\mathrm{I}_n$ & $v_\ell(n)$ & Steinberg$\otimes$unr. & $t=v_\ell(n)$, $e=0$ & Kim--Ota, $\ell\ge5$, (TW), $q\not\equiv\pm1$ or $\bar\rho$ ramified; data: all $\ell$, all $q$\\
$\mathrm{I}_n^*$ & $v_\ell(n)$ & $\chi_q\otimes$Steinberg & $v_\ell(n)+v_\ell(q^2-1)$ & Theorem~\ref{thm:euler} ($e$, all odd $\ell$), Theorem~\ref{thm:C} ($t$), Corollary~\ref{cor:both} (sum, $\ell\ge5$)\\
$\mathrm{I}_0^*$ & $0$ & $\chi_q\otimes$unramified & $v_\ell(q-1)+v_\ell((q+1)^2-a_q(E')^2)$ & Theorem~\ref{thm:euler}, all odd $\ell$\\
$\mathrm{IV},\mathrm{IV}^*$ & $[\ell=3]$ & tame, order $3$ & $v_\ell(q-\chi_{-3}(q))$ & qualitative: Theorem 1.5(a) of the note; quantitative: Proposition~\ref{prop:selmer}, data\\
$\mathrm{II},\mathrm{II}^*$ & $[\ell=3]$ & tame, order $6$ & $v_\ell(q-\chi_{-3}(q))$ & same as $\mathrm{IV}^*,\mathrm{IV}$ by Theorem~\ref{thm:twist}(c)\\
$\mathrm{III},\mathrm{III}^*$ & $0$ & tame, order $4$ & $v_\ell(q-\chi_{-4}(q))$ & attained in isolation; not additive in the configuration of Section~\ref{sec:anomaly}\\
$2,3$ & empirical & wild & PARI's factor; $t$ via $v_q(j)$ & data, no exception\\
$\ell$ & empirical & & $w_3$ & data\\
\bottomrule
\end{tabular}
\caption{Weights of the note, exact local terms and their status, by Kodaira type.}\label{tab:cmp}
\end{table}

\section{Changes made to the note (10 October 2026)}\label{sec:rec}
\begin{enumerate}
\item The former Conjectures 1.1 and 1.3 are replaced by Conjecture~\ref{conj:euler} (Conjecture 1.2 of the note); the component-group and twist-symmetrised forms are its Proposition 1.3.
\item Theorems~\ref{thm:twist}, \ref{thm:euler}, \ref{thm:C} and Corollary~\ref{cor:both} are Theorems 4.3, 4.4, 4.5 and Corollary 4.6 of the note, summarised as its Theorem 1.6; the former Theorem 1.4 (the qualitative statements at $\ell=3$) is its Theorem 1.5.
\item Section~\ref{sec:prop} is Section 5 of the note, with Proposition~\ref{prop:selmer} stated there as Conjecture 5.1 (Selmer form of the rule); the introduction presents the modular degree as governed by the naive adjoint $L$-value, the local terms as the omitted Euler factors, and the type $\mathrm{III}$ anomaly as the test case. This is the question for Kim and Ota, who have exactly the tools (their Theorem 2.13 and Proposition 2.15).
\item The equality rates of Table~\ref{tab:uniform} are reported in Section 6.3 of the note as the measure of sharpness.
\end{enumerate}

\begin{thebibliography}{9}
\bibitem{DDT} H.~Darmon, F.~Diamond, R.~Taylor, Fermat's last theorem, in: Current developments in mathematics, 1995, Int.\ Press, 1994, 1--154 (theorem numbers as in the revised version at \url{https://www.math.mcgill.ca/darmon/pub/Articles/Expository/05.DDT/paper.pdf}).
\bibitem{DFG04} F.~Diamond, M.~Flach, L.~Guo, The Tamagawa number conjecture of adjoint motives of modular forms, Ann.\ Sci.\ \'Ecole Norm.\ Sup.\ 37 (2004), 663--727.
\bibitem{KimOta} C.-H.~Kim, K.~Ota, On the quantitative variation of congruence ideals and integral periods of modular forms, Res.\ Math.\ Sci.\ 10 (2023), Paper No.\ 22 (arXiv:1905.02926).
\bibitem{RibetStein} K.~A.~Ribet, W.~A.~Stein, Lectures on Serre's conjectures, in: Arithmetic algebraic geometry (Park City, 1999), IAS/Park City Math.\ Ser.\ 9, Amer.\ Math.\ Soc., 2001, 143--232.
\end{thebibliography}
\end{document}
